\documentclass[a4paper,11pt]{article}

\usepackage[utf8]{inputenc}
\usepackage[T1]{fontenc}
\usepackage[italian]{babel}
\usepackage{geometry}
\usepackage{hyperref}
\usepackage{listings}
\usepackage{xcolor}
\usepackage{enumitem}
\usepackage{amsmath}

\geometry{margin=2.5cm}
\hypersetup{
    colorlinks=true,
    linkcolor=blue,
    urlcolor=blue,
    pdftitle={Relazione Progetto LSO -- Labirinto con oggetti e uscita}
}

\lstdefinestyle{cstyle}{
    language=C,
    basicstyle=\ttfamily\small,
    keywordstyle=\color{blue},
    commentstyle=\color{gray},
    stringstyle=\color{teal},
    breaklines=true,
    frame=single,
    columns=fullflexible,
    showstringspaces=false
}
\lstset{style=cstyle}

\title{\textbf{Relazione del progetto LSO}\\
Traccia A -- Labirinto con oggetti e uscita}
\author{
    Nomi componenti del gruppo \\
    \texttt{(inserire qui i nomi)}
}
\date{\today}

\begin{document}
\maketitle

\tableofcontents
\newpage

\section{Introduzione}

Il progetto implementa un sistema client--server concorrente in C per la gestione di
una partita in un labirinto 2D con oggetti, uscite e più giocatori contemporanei.
Il server mantiene la mappa completa, gestisce la posizione dei client, la scoperta
progressiva delle celle e il punteggio finale. Il client fornisce un'interfaccia
testuale interattiva (TUI) per login, registrazione, movimento, consultazione delle
mappe, elenco utenti e classifica.

La soluzione segue la traccia A, rispettando l'uso di socket TCP, la separazione
tra protocollo applicativo e livello di trasporto e la comunicazione asincrona
dal server verso i client. Il codice è organizzato in moduli separati: \texttt{server/}
(logica di gioco e gestione connessioni), \texttt{client/} (interfaccia utente e
parsing risposte) e \texttt{common/} (protocollo e costanti condivise).

\section{Guida d'uso}

\subsection{Compilazione}

Il progetto usa un \texttt{Makefile} con i seguenti target principali:

\begin{lstlisting}[language=bash]
make build     # compila client e server (debug)
make run-server # compila e avvia il server sulla porta 8080
make run-client # compila e avvia il client su 127.0.0.1:8080
make test      # compila ed esegue la suite di test (49 test)
\end{lstlisting}

La compilazione del server richiede il supporto ai thread POSIX (\texttt{-pthread})
e alla libreria \texttt{hiredis} per l'autenticazione via Redis. La compilazione
del client non richiede dipendenze esterne oltre alla libc standard.

Il server è composto da \texttt{server/server.c}, \texttt{server/game.c},
\texttt{server/log.c} e \texttt{common/protocol.c}. Il client da
\texttt{client/client.c}, \texttt{client/ui.c}, \texttt{client/command.c},
\texttt{client/response.c}, \texttt{client/state.c} e \texttt{common/protocol.c}.

\subsection{Avvio del server}

Il server accetta opzionalmente una porta e un percorso per il file di log.
I valori di default sono porta \texttt{8080} e log \texttt{logs/server.log}.

\begin{lstlisting}[language=bash]
./server_app                  # default: porta 8080
./server_app 9090             # porta personalizzata
./server_app 8080 mylog.txt   # porta e log personalizzati
\end{lstlisting}

All'avvio il server stampa a terminale:

\begin{lstlisting}[basicstyle=\ttfamily\small,frame=none]
=== LSO Server ===
Port:      8080
Log file:  logs/server.log
Redis:     connected to 127.0.0.1:6379
Server is running in LOBBY. Press Ctrl+C to stop.
\end{lstlisting}

Quindi:
\begin{itemize}[topsep=0.3em]
    \item inizializza il logger e il socket TCP di ascolto;
    \item si connette a Redis su \texttt{127.0.0.1:6379} per gestire le credenziali
          utente (se Redis non è disponibile, l'autenticazione viene disabilitata
          e il server stampa un avviso);
    \item avvia il thread periodico per l'invio delle mappe globali mascherate
          e il timer di sessione.
\end{itemize}

\subsection{Avvio del client}

\begin{lstlisting}[language=bash]
./client_app 127.0.0.1 8080
\end{lstlisting}

Il client usa \texttt{getaddrinfo()} per risolvere hostname e indirizzi numerici,
connettendosi via socket TCP. Subito dopo la connessione, il terminale viene messo
in modalità raw per consentire l'input immediato da tastiera (movimento diretto
con i tasti \texttt{W/A/S/D}).

\subsection{Uso interattivo}

Il client prevede due modalità commutabili con il tasto \texttt{TAB}:
\begin{itemize}[topsep=0.3em]
    \item \textbf{Command mode} -- i comandi vengono digitati e confermati con
          \texttt{Invio}. Viene mostrato un prompt \texttt{>} con l'input buffer
          corrente e una lista dei comandi disponibili.
    \item \textbf{Movement mode} -- i tasti \texttt{W/A/S/D} inviano immediatamente
          i movimenti al server. \texttt{G} richiede la mappa globale mascherata,
          \texttt{L} torna alla mappa locale. \texttt{Q} (maiuscola) invia \texttt{QUIT}
          e disconnette il client.
\end{itemize}

Comandi testuali supportati:
\begin{itemize}[topsep=0.3em]
    \item \texttt{register <nickname> <password>} -- registra un nuovo utente
    \item \texttt{login <nickname> <password>} -- autentica un utente esistente
    \item \texttt{ready} -- marca il giocatore come pronto nella lobby
    \item \texttt{start} -- (solo owner) avvia la partita
    \item \texttt{reset} -- (solo owner) riporta la sessione in lobby dopo la fine
    \item \texttt{rank} -- richiede la classifica
    \item \texttt{users} -- richiede la lista degli utenti connessi
    \item \texttt{local} -- richiede la mappa locale (finestra $5\times5$)
    \item \texttt{global} -- richiede la mappa globale mascherata
    \item \texttt{quit} -- disconnette il client
\end{itemize}

\section{Protocollo applicativo}

\subsection{Struttura generale}

Il protocollo è line-based: ogni messaggio è una riga testuale terminata da
\texttt{\textbackslash n}. I messaggi multi-linea (mappe, liste utenti, classifica)
terminano con una riga \texttt{END}. Prefissi testuali distinguono i tipi di risposta
(\texttt{OK}, \texttt{ERR}, \texttt{SESSION}, \texttt{MAP}, \texttt{USERS},
\texttt{RANK}, \texttt{TIME}).

\begin{table}[h]
\centering
\begin{tabular}{l|l|l}
\textbf{Prefisso} & \textbf{Esempio} & \textbf{Descrizione} \\
\hline
\texttt{OK} & \texttt{OK authenticated} & Operazione riuscita \\
\texttt{ERR} & \texttt{ERR wrong password} & Errore \\
\texttt{SESSION} & \texttt{SESSION STARTED} & Evento di sessione \\
\texttt{MAP LOCAL} & \texttt{MAP LOCAL 5 5} & Mappa locale (finestra $5\times5$) \\
\texttt{MAP GLOBAL} & \texttt{MAP GLOBAL 21 21} & Mappa globale mascherata \\
\texttt{USERS} & \texttt{USERS 3} & Lista utenti connessi \\
\texttt{RANK} & \texttt{RANK 3} & Classifica finale \\
\texttt{TIME} & \texttt{TIME 245} & Tempo rimanente in secondi \\
\end{tabular}
\caption{Prefissi del protocollo di risposta}
\end{table}

\subsection{Comandi client $\to$ server}

I comandi vengono tradotti dal modulo \texttt{client/command.c} nel formato
protocollare. Il client supporta anche varianti con slash iniziale
(\texttt{/ready} $\equiv$ \texttt{ready}) e gestisce il comando \texttt{help}
localmente senza inviare dati al server.

\subsection{Mappe locale e globale}

La mappa locale ($5\times5$ celle, centrata sul giocatore) mostra solo le
celle già scoperte dal giocatore. La mappa globale ($21\times21$ celle)
rappresenta l'intera matrice di gioco con le celle non ancora esplorate
sostituite dal simbolo \texttt{?} (cella nascosta). Il server invia la mappa
globale periodicamente ogni \texttt{T\_INTERVAL} (5 secondi) a tutti i client
durante la partita.

\subsection{Lista utenti e classifica}

La lista utenti è inviata automaticamente dopo login, register, ready, reset e
disconnessione. Ogni riga mostra il nickname con marcatori \texttt{[owner]} e
\texttt{[ready]}. La classifica combina i punteggi dei giocatori ancora connessi
e di quelli già disconnessi, ordinati per punteggio decrescente.

\section{Dettagli implementativi}

\subsection{Architettura concorrente del server}

Il server è multithread:
\begin{itemize}[topsep=0.3em]
    \item un thread per ogni client connesso (creato con \texttt{pthread\_create}
          e detached con \texttt{pthread\_detach});
    \item un thread periodico (\texttt{periodic\_global}) che ogni
          \texttt{T\_INTERVAL} secondi invia la mappa globale mascherata e il
          tempo rimanente a tutti i client;
    \item un thread \texttt{session\_timer} che decreta la fine della sessione
          dopo \texttt{SESSION\_DURATION} (300) secondi.
\end{itemize}

La sincronizzazione è realizzata con quattro mutex indipendenti:
\begin{itemize}[topsep=0.3em]
    \item \texttt{maze\_mutex} -- protegge la matrice del labirinto;
    \item \texttt{client\_mutex} -- protegge la lista dei client connessi;
    \item \texttt{score\_mutex} -- protegge la tabella dei punteggi;
    \item \texttt{session\_mutex} -- protegge lo stato della sessione
          e il timer.
\end{itemize}

La separazione dei mutex riduce la contesa tra i thread.

\subsection{Stati della sessione}

Il server gestisce tre stati:
\begin{itemize}[topsep=0.3em]
    \item \textbf{LOBBY} -- i client possono registrarsi, autenticarsi
          e dichiararsi pronti (\texttt{READY}). Solo l'owner può avviare
          la partita con \texttt{START}.
    \item \textbf{PLAYING} -- la partita è attiva. I client possono muoversi,
          raccogliere oggetti, consultare le mappe. Il server invia
          periodicamente la mappa globale e il tempo rimanente.
    \item \textbf{FINISHED} -- la partita è terminata. I client possono
          consultare la classifica (\texttt{RANK}). Solo l'owner può
          resettare la sessione (\texttt{RESET}) per tornare in LOBBY.
\end{itemize}

\subsection{Avvio della partita}

La partita parte quando l'owner invia \texttt{START} e \emph{tutti}
i giocatori non-owner connessi e autenticati hanno inviato \texttt{READY}.
Se qualche giocatore non è pronto, il server risponde con
\texttt{ERR not all players are ready}.

All'avvio:
\begin{enumerate}[topsep=0.3em]
    \item il labirinto viene generato (algoritmo ricorsivo di backtracking,
          matrice $21\times21$, percorsi larghi una cella tra pareti spesse);
    \item vengono piazzati 15 oggetti in celle vuote casuali e un'uscita
          su uno dei quattro bordi;
    \item a ogni client viene assegnata una posizione di spawn casuale
          (round-robin da un pool di celle walkable a coordinate dispari);
    \item viene avviato il timer di 300 secondi.
\end{enumerate}

\subsection{Generazione del labirinto}

Il labirinto è generato con backtracking ricorsivo (depth-first search):
\begin{enumerate}[topsep=0.3em]
    \item l'intera matrice viene riempita di muri (\texttt{\#});
    \item partendo dalla cella $(1,1)$, si sceglie una direzione casuale
          con passo 2 e si apre la parete intermedia;
    \item si ricorre sulla cella di destinazione;
    \item la randomizzazione delle direzioni a ogni passo garantisce
          labirinti diversi a ogni esecuzione.
\end{enumerate}

\subsection{Movimento, oggetti e uscita}

Il movimento è a singola cella nelle quattro direzioni cardinali. Il
server controlla che la cella di destinazione sia raggiungibile (non muro).
In caso di movimento valido:
\begin{itemize}[topsep=0.3em]
    \item l'area di raggio \texttt{VIEW\_RADIUS} (2) attorno al giocatore
          viene marcata come visibile in modo permanente;
    \item se la cella contiene un oggetto, questo viene raccolto e rimosso
          dalla mappa;
    \item se la cella è l'uscita, il giocatore viene marcato come
          \texttt{exit\_reached} (il suo marcatore scompare per gli altri
          giocatori e lui passa in modalità spettatore potendo osservare
          la mappa globale).
\end{itemize}

La sessione termina quando \emph{tutti} i giocatori hanno raggiunto l'uscita,
oppure allo scadere del timer (300 secondi), oppure se tutti i giocatori
si disconnettono.

\subsection{Classifica e punteggio}

Il punteggio è calcolato con la formula:
\[
\text{score} = \text{objects\_collected} \times 100 - \text{exit\_time\_seconds}
\]
dove \texttt{exit\_time\_seconds} è il tempo trascorso dall'inizio della
sessione al momento in cui il giocatore ha raggiunto l'uscita. I giocatori
che non raggiungono l'uscita non ricevono punteggio e sono ordinati
per numero di oggetti raccolti.

L'ordinamento usa \texttt{qsort} con comparatore personalizzato:
prima i giocatori usciti (per punteggio decrescente), poi i non usciti
(per numero di oggetti decrescente).

La risposta \texttt{RANK} include per ogni giocatore il numero di oggetti
raccolti, il tempo di uscita (se applicabile) e il punteggio calcolato.

\subsection{Autenticazione e registrazione}

Il server usa Redis come archivio esterno per le credenziali. La chiave
usata è \texttt{user:<nickname>} con valore la password in chiaro.

\begin{itemize}[topsep=0.3em]
    \item \textbf{Registrazione}: controlla che il nickname non sia già
          presente in Redis e che non inizi con \texttt{guest} (riservato
          ai client non autenticati). In caso positivo, memorizza la
          coppia nickname--password.
    \item \textbf{Login}: recupera la password da Redis e la confronta
          con quella fornita. Se il nickname non esiste risponde
          \texttt{ERR user not registered}; se la password non corrisponde
          risponde \texttt{ERR wrong password}.
    \item \textbf{Occupazione nickname}: controlla che nessun altro client
          connesso stia già usando quel nickname, per evitare conflitti.
\end{itemize}

Se Redis non è disponibile all'avvio, il server stampa un avviso e le
operazioni di autenticazione rispondono con un errore.

\subsection{Interfaccia utente (TUI)}

Il client disegna un'interfaccia testuale persistente usando sequenze
di escape ANSI e caratteri box-drawing UTF-8 (\texttt{┌─┐│└─┘}).
La schermata viene ridisegnata a ogni cambiamento di stato.

La TUI mostra:
\begin{itemize}[topsep=0.3em]
    \item in \textbf{LOBBY}: il titolo, le istruzioni per owner/non-owner,
          la lista dei giocatori connessi con colori distinti (owner in verde
          grassetto, giocatori pronti in giallo, guest in blu);
    \item in \textbf{PLAYING}: la mappa locale o globale con legenda colorata
          (\texttt{P}=giocatore in verde, \texttt{O}=altri in giallo,
          \texttt{*}=oggetti in ciano, \texttt{E}=uscita in magenta,
          \texttt{?}=nascosto in blu, \texttt{\#}=muro), il timer e
          i controlli disponibili;
    \item in \textbf{FINISHED}: la classifica con colori per le prime tre
          posizioni (oro, argento, bronzo) e il punteggio in magenta.
\end{itemize}

La funzione \texttt{clear\_screen()} invia \texttt{\textbackslash033[2J\textbackslash033[H}
per pulire lo schermo a ogni frame. La funzione \texttt{visible\_width()} gestisce
correttamente le sequenze multi-byte UTF-8 e i codici di escape ANSI per
mantenere l'allineamento dei bordi.

\subsection{Logging}

Il server registra in \texttt{logs/server.log} tutti gli eventi principali:
avvio, connessioni (\texttt{CONNECTED: <nick>}), disconnessioni
(\texttt{DISCONNECTED: <nick>}), inizio/fine/reset sessione.
Ogni riga è preceduta da un timestamp nel formato
\texttt{[YYYY-MM-DD HH:MM:SS]}. La scrittura è protetta da mutex
poiché più thread client possono scrivere contemporaneamente.
I messaggi di log vengono anche stampati su \texttt{stderr} per
facilitare il debug in tempo reale.

\subsection{Gestione delle disconnessioni}

Quando un client si disconnette durante la partita:
\begin{itemize}[topsep=0.3em]
    \item il suo punteggio viene salvato nella tabella \texttt{known\_scores}
          per apparire nella classifica finale;
    \item se era l'owner, la proprietà passa al primo client rimanente;
    \item se non rimangono client connessi, la sessione passa a
          \texttt{FINISHED};
    \item se tutti i giocatori rimanenti hanno già raggiunto l'uscita,
          la sessione termina immediatamente.
\end{itemize}

\section{Mappa dei moduli}

\begin{table}[h]
\centering
\begin{tabular}{l|l|p{6cm}}
\textbf{Modulo} & \textbf{File} & \textbf{Responsabilità} \\
\hline
Server core     & \texttt{server/server.c}   & Main loop, accept, gestione comandi, stati \\
Game logic      & \texttt{server/game.c}     & Labirinto, movimento, visibilità, mappe \\
Logging         & \texttt{server/log.c}      & Scrittura log con timestamp e mutex \\
\hline
Client core     & \texttt{client/client.c}   & Connessione, select loop, dispatch input \\
TUI             & \texttt{client/ui.c}       & Rendering schermo, bordi, mappe, colori \\
Commands        & \texttt{client/command.c}  & Traduzione comandi utente $\to$ protocollo \\
Response parser & \texttt{client/response.c} & Parsing risposte server multi-linea \\
Client state    & \texttt{client/state.c}    & Stato client, autenticazione pendente \\
\hline
Protocollo      & \texttt{common/protocol.h} & Costanti comandi/risposte, send\_all, recv\_line \\
Costanti        & \texttt{common/constants.h}& Dimensioni, simboli, colori ANSI \\
\end{tabular}
\caption{Organizzazione modulare del progetto}
\end{table}

\section{Scelte progettuali}

Le principali scelte progettuali sono state:
\begin{enumerate}[topsep=0.3em]
    \item \textbf{Separazione client/server/protocollo}: ogni livello è
          indipendente e testabile singolarmente. Il client non conosce
          la logica di gioco, il server non conosce l'interfaccia utente.
    \item \textbf{Protocollo testuale line-based}: semplifica il debugging
          (i messaggi sono leggibili con \texttt{telnet}) e il parsing
          (ogni riga è autoconclusiva fino a \texttt{\textbackslash n}).
    \item \textbf{Server multithread con mutex granulari}: un thread per
          client più thread ausiliari, con mutex separati per mappa,
          lista client, punteggi e stato sessione.
    \item \textbf{TUI senza ncurses}: l'interfaccia usa solo sequenze
          ANSI e caratteri UTF-8, senza dipendenze esterne, funzionante
          su qualsiasi terminale moderno.
    \item \textbf{Redis per credenziali}: archivio esterno leggero che
          persiste le registrazioni finché il server Redis è attivo.
          In caso di indisponibilità il server funziona comunque (con
          autenticazione disabilitata).
    \item \textbf{Test automatici}: suite di 49 test che coprono il parser
          comandi, lo stato client, il parser risposte, il protocollo e
          la logica di gioco (movimento, visibilità, pickup, uscita).
\end{enumerate}

\section{Conclusioni}

Il progetto soddisfa i requisiti della traccia A: gestione concorrente di più
client, labirinto con oggetti e uscita, scoperta progressiva delle celle,
invio periodico della mappa globale mascherata, autenticazione utenti,
logging e classifica finale con punteggio composito (oggetti $\times$ tempo
di uscita). La struttura modulare, la suite di test e l'uso di strumenti
standard (Makefile, Redis, ANSI escapes) rendono il sistema manutenibile
e portabile.

Tutto il codice sorgente, i test e la documentazione sono disponibili
nel repository del progetto.

\end{document}
